\documentclass[a4paper,11pt]{article}
\usepackage[margin=27mm]{geometry}
\usepackage{fontspec,amsmath,amssymb,amsthm,unicode-math,booktabs}
\setmainfont{STIXGeneral}\setmathfont{STIX Two Math}
\setlength{\emergencystretch}{2em}
\newcommand{\gen}[1]{\langle#1\rangle}
\newcommand{\Sym}{\operatorname{Sym}}
\newcommand{\Alt}{\operatorname{Alt}}
\newtheorem{theorem}{Theorem}\newtheorem{lemma}[theorem]{Lemma}
\newtheorem{corollary}[theorem]{Corollary}
\title{Rank augmentation from a supported transposition}
\author{Human Compression project}\date{Private draft; independent review pending}
\begin{document}\maketitle
\begin{abstract}
We study a string marking $(t,a,b,c)$ in which $t$ is a single
transposition and every orbit of $\gen{a,b}$ has at most three points.
The normal closure of $t$ gives the relevant end group explicitly.
A dihedral stabilizer argument proves that both ends have the intersection
property. Consequently the remaining obstruction is a small supported
kernel. This proves the four- and five-point cases of the rank augmentation
theorem of De Saedeleer--Leemans--Mulpas with their stated assumptions.
For the six-point case it gives an exact test involving seven subgroup memberships and a useful
sufficient orbit condition, replacing the false unrestricted assertion.
All conclusions here are written proofs; a finite example is not used to
establish a quantified statement.
\end{abstract}

\section{The common setting}
Let $\Omega$ be finite and let $t,a,b,c$ be distinct nonidentity
involutions in $\Sym(\Omega)$. Suppose that
\[
 t=(\xi\,\eta),\qquad b(\xi)=\xi,\ b(\eta)=\eta,\ c(\xi)=\xi,\ c(\eta)=\eta,
 \qquad ac=ca,
\]
and that $G=\gen{t,a,b,c}$ is transitive. Put
\[
 B=\gen{a,b},\qquad H=\gen{t,a,b},\qquad K=\gen{a,b,c}.
\]
Assume that every $B$-orbit has size at most three. Let $O$ be the
$H$-orbit containing $\xi,\eta$. We first show that its size is four,
five or six; this is a consequence, not an additional assumption.
All permutation products below act on the right.

The components for the two involutions $a,b$ are singletons, one edge,
a double edge, or a three-vertex path with successive labels $a,b$.
At a vertex fixed by $b$, its component is a singleton, an $a$-edge,
or the end of such a path. Each $a$-neighbor of an endpoint of $t$
is fixed by $c$, because $a$ commutes with $c$. The orbit $O$ is precisely
the union of the $B$-components containing the two endpoints. If neither
component were a triple, $c$ would fix this union pointwise and $O$ would
be a $G$-orbit. Transitivity would give $O=\Omega$ and $c=1$, a contradiction.
Thus at least one component is a triple. Such a component has only one
point fixed by $b$, so the two endpoints of $t$ lie in distinct
$B$-components. Their union has size four, five or six. The marked paths
on $O$, after a possible reversal, are exactly
\[
\begin{array}{ccl}
4&:&x\mathrel{-_t}u\mathrel{-_a}v\mathrel{-_b}w,\\
5&:&x\mathrel{-_a}y\mathrel{-_t}u\mathrel{-_a}v\mathrel{-_b}w,\\
6&:&\ell\mathrel{-_b}k\mathrel{-_a}y\mathrel{-_t}u
                   \mathrel{-_a}v\mathrel{-_b}w.
\end{array}                                                    \tag{1}
\]
The letters in (1) are local point names. In each row, $T=\{u,v,w\}$ is a triple with
$u$ fixed by $b,c$ and $v$ fixed by $c$. In the six-point row the left
triple has the same property with its order reversed.

We use the classical end criterion for string C-groups: two C-marked
ends whose intersection is their common marked subgroup give a C-marking
of the whole group. This is McMullen--Schulte, Proposition~2E16(a),
quoted in [DLM, Proposition~5.1]. A pair of distinct involutions is a
rank-two C-marking. We retain this criterion, rather than presenting it
as new.

\section{The two ends}
Every three-point component gives the same marked quotient $S_3$ of $B$.
On a two-point component, $ab$ has order at most two. Consequently $ab$
has order three or six, and $B$ is dihedral of order six or twelve.
Write
\[
 z=(ab)^3.
\]
Then $z$ is either the identity or a central involution of $B$; it fixes
every triple pointwise. The kernel of the action of $B$ on any one triple
is $\gen z$. In particular, the stabilizers in $B$ of $u$ and $w$ are
$\gen{b,z}$ and $\gen{a,z}$ respectively.

\begin{lemma}\label{reflection}
Let $D=\gen{p,q}$ be generated by two involutions. Suppose $q$ fixes
$v$, while $p$ sends $v$ to $w\ne v$ and $q$ does not fix $w$.
An element of $D$ fixing both $v,w$ belongs to its cyclic rotation
subgroup.
\end{lemma}
\begin{proof}
Put $r=qp$, so $v^r=w$. Since $q$ fixes $v$ and normalizes $\gen r$,
the $D$-orbit of $v$ is its $\gen r$-orbit. This orbit has at least three
points: it contains $v,w,w^q$, and $w^q$ can equal neither $w$ nor $v$
(the latter equality would contradict $v^q=v$). Every element outside
$\gen r$ conjugates $r$ to $r^{-1}$. If such an element $d$ fixed both
$v,w$, then
\[
 v^r=w=w^d=(v^r)^d=(v^d)^{r^{-1}}=v^{r^{-1}}.
\]
It would follow that $v^{r^2}=v$, contrary to the orbit size.
\end{proof}

\begin{lemma}\label{ends}
Under the common hypotheses, both $(a,b,c)$ and $(t,a,b)$ are
string C-markings.
\end{lemma}
\begin{proof}
First consider $B\cap\gen{b,c}$. Every element of $\gen{b,c}$ fixes
$u$, so the intersection lies in $\gen{b,z}$. If it were larger than
$\gen b$, it would contain a nonidentity $z$.
In the five-point row, $b,c$ fix the left pair $\{x,y\}$ pointwise,
whereas $z=(ab)^3$ interchanges it. This already gives a contradiction.

For the four- and six-point rows, choose an anchored triple whose terminal
point is moved by $c$. Such a triple exists: otherwise $c$ fixes every
point of $O$, making $O$ a $G$-orbit and forcing $c=1$. Apply
Lemma~\ref{reflection} to $p=b,q=c$ and the last two points $v,w$ of
that triple. Since $z$ fixes both, it is a rotation in $\gen{b,c}$.
It commutes with $b$, since it is central in $B$, and hence also with
$c$. It already commutes with $a,b$. Moreover $z$ fixes the endpoints
of $t$ in these two rows, so it commutes with $t$ as well. A central
permutation of a transitive group that fixes one point fixes every point:
if $h(v)=v$, then $h(g(v))=g(h(v))=g(v)$. Thus $z=1$, again a contradiction.
We have proved $B\cap\gen{b,c}=\gen b$.

Next consider $B\cap\gen{t,a}$. This subgroup fixes $w$, so it lies
in $\gen{a,z}$. A larger intersection than $\gen a$ would again put
nonidentity $z$ in the intersection. Apply Lemma~\ref{reflection} with
$p=a,q=t$ at $v,u$. The $t,a$-orbit containing them has size three in
the four-point row and size four in the other two rows. Hence $z$ is a
rotation in $\gen{t,a}$. On this orbit the rotation is a full three-cycle
or four-cycle. Since $z$ fixes $u$, its exponent is respectively divisible
by three or four. Outside this orbit, $t$ fixes
every point, and the rotation acts as $a$, of order at most two.
In the five- and six-point rows, divisibility by four therefore gives
$z=1$ globally. In the four-point row an even exponent also gives $z=1$;
the only other possibility is the restriction $a_{\mathrm{out}}$ of $a$
to the complement of the three-point $t,a$-orbit.

In this last possibility, $a_{\mathrm{out}}=a(u\,v)$. Since $c$ fixes
$u,v$ and commutes with $a$, it commutes with $a_{\mathrm{out}}$.
The permutation $a_{\mathrm{out}}=z$ commutes with $a,b$, and its support
is disjoint from that of $t$. It is again central in $G$ and has a fixed
point, so it is the identity. Thus $B\cap\gen{t,a}=\gen a$.
Both ends now satisfy the retained rank-three end criterion.
\end{proof}

\section{The supported factor}
If a permutation is defined on a subset of $\Omega$, we extend it by the
identity on the complement. Let $B_{\mathrm{out}}$ be the restriction image
of $B$ on $\Omega\setminus O$, viewed in this way as a permutation group.

\begin{lemma}\label{factor}
In the four- and five-point rows,
\[
 H=\Sym(O)\times B_{\mathrm{out}}.                         \tag{2}
\]
In the six-point row let $f$ be the unique bijection from the left triple
to the right triple preserving the marked $a,b$ actions. Let
\[
 N=\gen{(v\ f(v)):v\hbox{ in the left triple}}\cong C_2^3.
\]
Then
\[
 H=N\rtimes B.                                             \tag{3}
\]
These are identities of the actual permutation subgroups, not only
abstract group isomorphisms.
\end{lemma}
\begin{proof}
Conjugating the pure transposition $t$ by $B$ gives transpositions on
its orbit of unordered pairs. In the four-point row these pairs are the
three edges of the star from $x$ to $T$. They generate $\Sym(O)$.
In the five-point row, the action of $B$ on the left pair and the right
triple is the full product $C_2\times S_3$. Indeed $z$ interchanges only
the left pair on $O$, and the quotient action on $T$ is $S_3$.
The conjugate pairs are therefore all six edges of the complete bipartite
graph between the pair and the triple. Again their transpositions generate
$\Sym(O)$. In both cases these transpositions fix the complement of $O$
pointwise. For any $g\in B$, multiply by the inverse of its restriction
to $O$, already available in $\Sym(O)$, to obtain its pure complementary
restriction. This proves (2), including independence of the two factors.

In the six-point row, $B$ acts on its two triples through the same
marked $S_3$ action, identified by $f$. The conjugates of $t$ are exactly
the three disjoint matching transpositions $(v\ f(v))$. They generate
$N$, and $B$ permutes these three generators. No nonidentity member of
$N$ preserves the two triples separately, whereas every member of $B$
does. Thus $N\cap B=1$ and (3) follows from $H=\gen{t,B}$.
\end{proof}

The remaining intersection can be expressed without testing all
parabolic pairs. In (2), every $h\in H\cap K$ has the same complementary
restriction as some $g\in B$. Hence $hg^{-1}$ is supported on $O$.
It follows that
\[
 H\cap K=B\quad\Longleftrightarrow\quad
 K\cap\Sym(O)=B\cap\Sym(O).                                \tag{4}
\]
This is the elementary restriction-kernel test: the actual images are
already known, so only their kernels remain to be compared. In (3), the
unique factorization $h=ng$ gives instead
\[
 H\cap K=(N\cap K)B,
 \qquad H\cap K=B\quad\Longleftrightarrow\quad N\cap K=1.   \tag{5}
\]
The product in (5) is a semidirect product, since $B$ normalizes $N\cap K$.

\section{Blocks and parity in the four- and five-point rows}
Removing the single $t$-edge disconnects the small left component in
(1), a singleton or a pair. In either row $K$ fixes that component setwise.
Every other point belongs to one $K$-orbit, say $R$: the full labelled
graph is connected, and the removed edge was its sole edge at that small
component. In particular every triple outside $O$ lies in $R$.

\begin{lemma}\label{blocks}
Let $R$ be any $K$-orbit containing an anchored triple $T$, and suppose
$K$ acts imprimitively on $R$. Every $B$-triple in $R$ has its three
points in distinct blocks of each nontrivial system of blocks.
Moreover, there is another $B$-triple $T'\subseteq R\setminus T$
whose three points lie respectively in the same three blocks as $T$.
In the four- and five-point rows, $T'$ lies outside $O$.
\end{lemma}
\begin{proof}
The equivalence relation induced on any $B$-triple by the blocks is
invariant under its full $S_3$ action. It is either universal or discrete.
Suppose a triple $U$ lay in one block $D$. Then $a,b$ preserve $D$.
If $c$ also preserved $D$, transitivity on $R$ would give $D=R$.
Otherwise put $D'=D^c\ne D$. These blocks are disjoint, $c$ interchanges
them, and $a$ preserves $D'$ because it commutes with $c$.
The images under $c$ of the two $a$-adjacent points of $U$ are
$a$-adjacent points $v',w'$ in $D'$. They cannot both be sent outside
$D'$ by $b$: their two distinct images, together with $v',w'$, would
belong to one $B$-orbit of at least four points. Thus $b$ sends a point
of $D'$ into $D'$, and hence preserves that whole block. The union
$D\cup D'$ is now invariant under $a,b,c$, so equals $R$.
But $c$ interchanges its two blocks, contradicting the fixed point $u$
of $c$ in the anchored triple. This proves the first assertion.

All blocks have the same size, at least two. Choose $u'\ne u$ in the
block of $u$. The points $u'$, $v'=(u')^a$, and $w'=(u')^{ab}$ lie in
the three distinct blocks of $u,v,w$. They belong to one $B$-orbit, which
has at most three points. Its marked action is therefore the triple
$u'\mathrel{-_a}v'\mathrel{-_b}w'$. This orbit is disjoint from $T$:
if it met $T$ the two $B$-orbits would coincide, whereas $u'$ is not any
point of $T$. In the four- and five-point rows the other points of $O$
belong to the small left $K$-orbit, so this new triple is outside $O$.
\end{proof}

For comparison with [DLM, Lemma~5.3], every imprimitive $K$-orbit in
our setting contains an anchored triple. Indeed, removing one edge from
the connected labelled graph leaves at most two components, each
containing an endpoint of that edge. If its $B$-component is a singleton
or a pair, $c$ fixes that whole component and its $K$-orbit has that same
size, admitting no nontrivial block system. Otherwise the component is
an anchored triple. Thus Lemma~\ref{blocks} proves both conclusions of
the source lemma without its extended graph propagation argument.

\begin{theorem}\label{positive}
In the common setting, the output $(t,a,b,c)$ is a string C-marking if
one of the following holds.
\begin{enumerate}
\item $|O|=4$, and $K$ acts imprimitively on its orbit containing $T$
whenever that orbit contains another $B$-triple.
\item $|O|=5$, the permutations $a,b,c$ are even, and the same
imprimitivity condition holds.
\end{enumerate}
\end{theorem}
\begin{proof}
By Lemma~\ref{ends}, it suffices to prove $H\cap K=B$.
Suppose first that there is a triple outside $O$. The relevant $K$-orbit
is $R$, so the assumed imprimitivity and Lemma~\ref{blocks} give a triple
$T'$ outside $O$ sharing the three blocks of $T$. Take $h\in H\cap K$,
and use (2) to choose $g\in B$ with the same action outside $O$.
Then $h'=hg^{-1}$ fixes $T'$ pointwise and lies in $K$; it consequently
preserves all three specified blocks. It also preserves $T$ setwise,
since it preserves $O$ and the small left $K$-orbit. Each of these blocks
contains exactly one point of $T$, so $h'$ fixes $T$ pointwise.
In the four-point row the remaining point of $O$ is fixed by $K$, whence
$h'=1$. In the five-point row the only other possibility is the single
transposition of the left pair. This is excluded by $K\leq\Alt(\Omega)$.
Thus $h=g\in B$.

It remains to handle the case with no triple outside $O$, so no
imprimitivity assumption is available. Let $a_{\mathrm{out}}$ be the
restriction of $a$ to $\Omega\setminus O$. On that complement all
$B$-orbits have size at most two; hence $a_{\mathrm{out}}$ commutes with
$b$. It commutes with $a$ and $t$ as well. The restriction of $a$ to $O$
has support fixed pointwise by $c$, so $ac=ca$ also gives
$a_{\mathrm{out}}c=ca_{\mathrm{out}}$. Centrality in transitive $G$ and
its fixed points on $O$ force $a_{\mathrm{out}}=1$. Thus
$B_{\mathrm{out}}=\gen{b_{\mathrm{out}}}$ has order at most two.

For $|O|=4$, $B=\Sym(T)\times\gen{b_{\mathrm{out}}}$.
Indeed, if $b_{\mathrm{out}}\ne1$, then $(ab)^3=b_{\mathrm{out}}$
isolates its complementary factor; the restrictions of $a,b$ generate
$\Sym(T)$. Every element of $H\cap K$ fixes the remaining point of $O$
and has complementary action in $\gen{b_{\mathrm{out}}}$, so belongs to
this very group $B$.

For $|O|=5$, evenness of $b$ forces $b_{\mathrm{out}}\ne1$, since its
restriction to $O$ is a single transposition. The group
\[
 L=\Sym(\{x,y\})\times\Sym(T)\times\gen{b_{\mathrm{out}}}
\]
has order $24$. Its even subgroup has order $12$, and equals $B$:
$B$ has order $12$ by its faithful $C_2\times S_3$ action on $O$, and
both $a,b$ are even. Every member of $H\cap K$ preserves the left pair
and $T$, has complementary action in $\gen{b_{\mathrm{out}}}$, and is
even. It therefore belongs to $B$. This proves both cases, including the
boundary where the source's imprimitivity condition is vacuous.
\end{proof}

\section{An exact replacement for the six-point assertion}
\begin{theorem}\label{six}
In the six-point row of (1), the marking $(t,a,b,c)$ has the
intersection property if and only if none of the seven nonidentity
products of the three matching transpositions belongs to $K$.
In particular it has that property if the left and right triples of $O$
lie in distinct $K$-orbits.
\end{theorem}
\begin{proof}
Both ends have the intersection property by Lemma~\ref{ends}.
Equation (5) and the end criterion show that the whole marking has it
exactly when $N\cap K=1$. Since $N$ is generated by three disjoint
transpositions, its seven nonidentity elements are precisely the listed
products. If the two triples lie in different $K$-orbits, each nonidentity
member of $N$ moves a point from one such orbit to the other and so cannot
belong to $K$.
\end{proof}

The last condition is a genuine sufficient restriction, not a consequence
of the source hypotheses. The separately frozen eighteen-point example
has $K=S_{18}$ and contains $t$ itself; it fails the first of the seven
tests. Nor does the theorem above merely repeat the full intersection
property: it replaces all parabolic intersections by membership in one
explicit group of order eight, after separately proving the two end
properties uniformly.

\section{The source application and attribution}
Start with the exact hypotheses of [DLM, Theorem~4.1]: a natural
$S_n$ rank-three C-marking $(a,q,c)$, $n\geq5$, with $q$ not a
transposition. Choose a transposition $t$ in $q$ whose endpoints are
fixed by $c$, and put $b=tq$. Then $b$ is an involution, fixes those
endpoints, and commutes with $t$; also $c$ commutes with $t$ and $a$.
The new marking generates $S_n$, since $q=tb$. The orbit argument in
Section~1 proves, without selecting one of the three cases, that the
orbit containing $t$ has size four, five or six. The local paths (1)
also show that the four marks are distinct and nonidentity. Thus all
common hypotheses above are satisfied when every $B$-orbit has size at
most three. Theorem~\ref{positive} proves the source's cases (1)
and (2) with their full original assumptions. Theorem~\ref{six} gives the
stated corrected replacement for case (3), whose unrestricted assertion
is refuted in the separate exact counterexample.

The rank augmentation operation, cases (1) and (2), local component
analysis and intended end-intersection method belong to De Saedeleer,
Leemans and Mulpas. Our normal-closure and restriction-kernel formulations
reuse the elementary supported-transposition method from our common
gluing proof. The short dihedral and block arguments here replace the
longer local propagation arguments needed for these conclusions; we do
not claim that the underlying group facts are historically new.
The standard end criterion remains an attributed external premise.
Two declarations in \texttt{SupportedIntersections.lean} prove the actual
arbitrary-group factorization tests behind (4) and (5), using subgroups
and their intersections. The actual permutation factorizations in
Lemma~\ref{factor}, the end arguments, block and parity arguments, and
source application are written proofs; they are not silently included
in that formal certificate.
The proof includes both conclusions of source Lemma~5.3. The companion application verifies the source's three Section~6 examples
using its complete required program. The introductory Table~1 remains
explicitly retained motivating data; it is not a newly regenerated census.

\begin{thebibliography}{9}
\bibitem{DLM} J. De Saedeleer, D. Leemans and J. Mulpas,
\emph{A rank augmentation theorem for rank 3 string C-group representations
of the symmetric groups}, arXiv:2212.12563v1, 23 December 2022,
Theorem~4.1 and Section~5.
\end{thebibliography}
\end{document}
